\section{Deployment}

Questa sezione descrive come il sistema viene eseguito, configurato e avviato. L'obiettivo e' mostrare che l'architettura progettata puo' essere effettivamente messa in funzione in un ambiente coerente con le scelte tecnologiche effettuate.

\subsection{Architettura di esecuzione}

KompozeR prevede l'esecuzione coordinata di frontend, gateway, microservizi backend e basi di dati dedicate. In questa sottosezione e' utile descrivere come tali componenti vengano avviati e collegati tra loro in ambiente locale o di test.

Qui si puo' inserire un diagramma di deployment semplificato che mostri container o processi principali e relative dipendenze.

\subsection{Containerizzazione dei servizi}

Se il progetto utilizza Docker, questa parte dovrebbe spiegare come ogni componente applicativo venga pacchettizzato e reso eseguibile in modo consistente. Conviene chiarire:

\begin{itemize}
	\item quali componenti vengono containerizzati;
	\item il ruolo delle immagini applicative;
	\item il vantaggio in termini di isolamento e riproducibilita'.
\end{itemize}

\subsection{Configurazione dell'ambiente}

Per l'esecuzione del sistema e' necessario definire variabili d'ambiente, porte di ascolto, URL dei servizi e parametri di connessione ai database. In questa sottosezione conviene riassumere la configurazione minima richiesta, evitando di riportare dati sensibili ma chiarendo la logica generale del setup.

\subsection{Avvio del sistema in locale}

Questa parte dovrebbe descrivere in modo sintetico la procedura necessaria per eseguire il progetto in ambiente di sviluppo o dimostrazione. Ad esempio:

\begin{itemize}
	\item preparazione dell'ambiente;
	\item avvio dei servizi e dei database;
	\item avvio del frontend;
	\item verifica del corretto funzionamento iniziale.
\end{itemize}

Se utile, qui si possono richiamare i principali comandi o i file di configurazione che rendono possibile l'avvio coordinato del sistema.

\subsection{Persistenza e dati iniziali}

In questa sottosezione si puo' spiegare come vengono gestiti i dati minimi necessari all'avvio del sistema, ad esempio utenti amministrativi iniziali, prodotti di catalogo di esempio o configurazioni di base utili per test e dimostrazione.

\subsection{Considerazioni operative}

La sezione puo' chiudersi con alcune osservazioni sulle principali implicazioni operative della soluzione adottata: facilita' di setup, modularita' dell'ambiente, possibilita' di sostituzione o aggiornamento dei singoli servizi e supporto a una futura evoluzione del sistema.